\subsection{半单模的描述}

在本节余下部分，A 是一个环，S 是单 A-模的类组成的集合。对每个 \(\lambda \in S\)，我们选取一个单模 \(S_{\lambda}\)，其类为 \(\lambda\)（例如 \(S_{\lambda} = \lambda\)）；把 \(S_{\lambda}\) 的自同态体之反环记作 \(D_{\lambda}\)。我们把 \(S_{\lambda}\) 视为 \((A, D_{\lambda})\)-双模。

设 \(M\) 是一个 A-模。对每个 \(\lambda \in \mathcal{S}\)，\(\mathrm{Hom}_{\mathrm{A}}(\mathrm{S}_{\lambda}, M)\) 是体 \(D_{\lambda}\) 上的左向量空间。由 VIII, 第 59 页及 II, §1, No.~6, 第 202 页的命题 6，存在唯一的 A-线性映射，称为典范映射，

\[
\alpha_ {\mathrm{M}}: \bigoplus_ {\lambda \in \mathcal {S}} \left(\mathrm{S} _ {\lambda} \otimes_ {\mathrm{D} _ {\lambda}} \operatorname{Hom} _ {\mathrm{A}} (\mathrm{S} _ {\lambda}, \mathrm{M})\right) \to \mathrm{M}
\]

它满足关系式

\[
\alpha_ {\mathrm{M}} (s \otimes f) = f (s)\tag{10}
\]

其中 \(\lambda \in \mathcal{S}\)，\(s \in S_{\lambda}\)，\(f \in \mathrm{Hom}_{\mathrm{A}}(\mathrm{S}_{\lambda}, \mathrm{M})\)。若给 \(\bigoplus_{\lambda \in \mathcal{S}} (\mathrm{S}_{\lambda} \otimes_{\mathrm{D}_{\lambda}} \mathrm{Hom}_{\mathrm{A}}(\mathrm{S}_{\lambda}, \mathrm{M}))\) 与 \(\mathrm{M}\) 赋以它们自然的 \(\mathrm{End}_{\mathrm{A}}(\mathrm{M})\)-模结构，则映射 \(\alpha_{\mathrm{M}}\) 是 \(\mathrm{End}_{\mathrm{A}}(\mathrm{M})\)-线性的。

命题 7. —— 设 M 是一个 A-模。典范映射 \(\alpha_{\mathrm{M}}\) 是单射。对每个 \(\lambda \in \mathcal{S}\)，在 \(\alpha_{\mathrm{M}}\) 下，\(\mathrm{S}_{\lambda} \otimes_{\mathrm{D}_{\lambda}} \mathrm{Hom}_{\mathrm{A}}(\mathrm{S}_{\lambda}, \mathrm{M})\) 的像是 M 的 \(\lambda\) 型同型分量。\(\alpha_{\mathrm{M}}\) 的像是 M 的底座。A-模 M 是半单的当且仅当映射 \(\alpha_{\mathrm{M}}\) 是双射。

设 \(\lambda \in \mathcal{S}\)。把 M 的 \(\lambda\) 型同型分量记作 \(\mathrm{M}_{\lambda}\)。从 \(\mathrm{S}_{\lambda}\) 到 M 的每个 A-线性映射都取值于 \(\mathrm{M}_{\lambda}\)（VIII, 第 66 页，命题 5）。于是，由 VIII, 第 62 页的命题 3, a)，映射 \(\alpha_{\mathrm{M}}\) 诱导出从 \(\mathrm{S}_{\lambda} \otimes_{\mathrm{D}_{\lambda}} \mathrm{Hom}_{\mathrm{A}}(\mathrm{S}_{\lambda}, \mathrm{M})\) 到 \(\mathrm{M}_{\lambda}\) 的双射。由于 M 的底座是族 \((\mathrm{M}_{\lambda})_{\lambda \in \mathcal{S}}\) 的直和，且模 M 是半单的当且仅当它等于自身的底座，命题得证。

定义 5. —— 设 M 是一个半单 A-模。M 的（关于族 \((\mathrm{S}_{\lambda})_{\lambda\in\mathcal{S}}\) 的）描述是一个偶 \(((\mathrm{V}_{\lambda})_{\lambda\in\mathcal{S}},\alpha)\)，其中 \(V_{\lambda}\) 是体 \(D_{\lambda}\) 上的左向量空间（对每个 \(\lambda\in S\)），而 \(\alpha:\bigoplus_{\lambda\in\mathcal{S}}(\mathrm{S}_{\lambda}\otimes_{\mathrm{D}_{\lambda}}\mathrm{V}_{\lambda})\to\mathrm{M}\) 是 A-模的同构。

由命题 7，每个半单 A-模 M 都有一个典范描述：它是偶 \(\left(\left(\mathrm{Hom}_{\mathrm{A}}\left(\mathrm{S}_{\lambda},\mathrm{M}\right)\right)_{\lambda \in \mathcal{S}},\alpha_{\mathrm{M}}\right)\)，其中 \(\alpha_{\mathrm{M}}\) 是由公式 (10) 定义的 A-线性映射。

命题 8. —— 设 M 是一个半单 A-模，\(((\mathrm{V}_{\lambda})_{\lambda \in \mathcal{S}}, \alpha)\) 是 M 的一个描述。

\begin{enumerate}
\def\labelenumi{\alph{enumi})}
\item
  对每个 \(\lambda \in \mathcal{S}\)，映射 \(\alpha\) 诱导出从 A-模 \(S_{\lambda} \otimes_{D_{\lambda}} V_{\lambda}\) 到 M 的 \(\lambda\) 型同型分量的同构。
\item
  对每个 \(\lambda \in \mathcal{S}\)，映射 \(\beta_{\lambda} : \mathrm{V}_{\lambda} \to \mathrm{Hom}_{\mathrm{A}}(\mathrm{S}_{\lambda}, \mathrm{M})\)（由 \(\beta_{\lambda}(v)(s) = \alpha(s \otimes v)\) 定义）是 \(\mathrm{D}_{\lambda}\)-线性的双射。
\item
  设 \(\mathbf{N}\) 是 \(\mathbf{M}\) 的一个子模。存在唯一的族 \((\mathrm{W}_{\lambda})_{\lambda \in \mathcal{S}}\) 具有下列性质：\(\mathrm{W}_{\lambda}\) 是一个 \(\mathrm{D}_{\lambda}\)-线性子空间，含于 \(\mathrm{V}_{\lambda}\)（对每个 \(\lambda \in \mathcal{S}\)），且 \(\mathbf{N}\) 是 \(\alpha\) 作用于模 \(\bigoplus_{\lambda \in \mathcal{S}} (\mathrm{S}_{\lambda} \otimes_{\mathrm{D}_{\lambda}} \mathrm{W}_{\lambda})\) 所得的像，这里该模与其在模 \(\bigoplus_{\lambda \in \mathcal{S}} (\mathrm{S}_{\lambda} \otimes_{\mathrm{D}_{\lambda}} \mathrm{V}_{\lambda})\) 中的典范像等同。对每个 \(\lambda \in \mathcal{S}\)，\(\mathrm{W}_{\lambda}\) 是这样的元素 \(v \in \mathrm{V}_{\lambda}\) 组成的集合：\(\alpha(s \otimes v)\) 属于 \(\mathbf{N}\)（对每个 \(s \in \mathrm{S}_{\lambda}\)）。
\end{enumerate}

A-模 \(S_{\lambda} \otimes_{D_{\lambda}} V_{\lambda}\) 都是 \(\lambda\) 型同型模（对每个 \(\lambda \in S\)）。于是断言 a) 由 \(\alpha\) 是同构以及 M 是族 \((M_{\lambda})_{\lambda \in \mathcal{S}}\) 的直和（VIII, 第 65 页，命题 4, a)）这两个事实得出。

设 \(\lambda \in \mathcal{S}\)。把 \(\alpha\) 在 \(S_{\lambda} \otimes_{D_{\lambda}} V_{\lambda}\) 上的限制记作 \(\alpha_{\lambda}: S_{\lambda} \otimes_{D_{\lambda}} V_{\lambda} \to M\)。把 No.~3 的记号应用于 \((A, D_{\lambda})\)-双模 \(S_{\lambda}\)，我们有

\[
\beta_ {\lambda} = \gamma (\alpha_ {\lambda}) = \mathcal {H} (\alpha_ {\lambda}) \circ \beta_ {\mathrm{V} _ {\lambda}},
\]

其中最后一式由 VIII, 第 60 页得出。于是，\(\beta_{\lambda}\) 是 \(D_{\lambda}\)-线性同态 \(\mathcal{H}(\alpha_{\lambda})\) 与 \(\beta_{V_{\lambda}}\) 的复合。于是断言 b) 由 VIII, 第 62 页的命题 3, b) 得出。

设 \(N\) 是 \(M\) 的一个子模。我们有 \(N = \bigoplus_{\lambda \in \mathcal{S}} (N \cap M_{\lambda})\)（VIII, 第 66 页，推论），故 c) 由 VIII, 第 62 页的定理 2 得出。

推论. —— 设 M 是一个半单 A-模。对 M 的每个子模 N 及 S 的每个元素 \(\lambda\)，把 \(\operatorname{Hom}_{\mathrm{A}}(\mathrm{S}_{\lambda},\mathrm{N})\) 与一个 \(D_{\lambda}\)-线性子空间等同起来，该子空间由 \(\operatorname{Hom}_{\mathrm{A}}(\mathrm{S}_{\lambda},\mathrm{M})\) 中像含于 N 的映射组成。

\begin{enumerate}
\def\labelenumi{\alph{enumi})}
\item
  映射 \(\mathrm{N}\mapsto(\mathrm{Hom}_{\mathrm{A}}(\mathrm{S}_{\lambda},\mathrm{N}))_{\lambda\in\mathcal{S}}\) 是从 M 的 A-子模组成的集合到满足下列条件的族 \((\mathrm{W}_{\lambda})_{\lambda\in\mathcal{S}}\) 组成的集合的双射：对每个 \(\lambda\in S\)，\(W_{\lambda}\) 是一个 \(D_{\lambda}\)-线性子空间，含于 \(\mathrm{Hom}_{\mathrm{A}}(\mathrm{S}_{\lambda},\mathrm{M})\)。
\item
  逆双射把族 \((\mathrm{W}_{\lambda})_{\lambda \in \mathcal{S}}\) 映到 M 的 A-子模 \(\sum_{\lambda \in \mathcal{S}}\sum_{w\in \mathrm{W}_{\lambda}}w(\mathrm{S}_{\lambda})\)。
\end{enumerate}

这是把命题 8, c) 应用于 M 的典范描述的一个重新表述。

命题 9. —— 设 M 与 M' 是半单 A-模，并分别考虑 M 与 M' 的描述 \(((\mathrm{V}_{\lambda})_{\lambda\in\mathcal{S}},\alpha)\) 与 \(((\mathrm{V}_{\lambda}^{\prime})_{\lambda\in\mathcal{S}},\alpha^{\prime})\)。对每个族 \(\boldsymbol{f}=(f_{\lambda})_{\lambda\in\mathcal{S}}\)（取自 \(\prod_{\lambda\in\mathcal{S}}\mathrm{Hom}_{\mathrm{D}_{\lambda}}(\mathrm{V}_{\lambda},\mathrm{V}_{\lambda}^{\prime})\)），存在唯一的 A-线性映射 \(\varphi(\boldsymbol{f})\in\mathrm{Hom}_{\mathrm{A}}(\mathrm{M},\mathrm{M}^{\prime})\) 使下面的图式交换：

\[
\begin{array}{c} \bigoplus_ {\lambda \in \mathscr {S}} (S _ {\lambda} \otimes_ {D _ {\lambda}} V _ {\lambda}) \xrightarrow {\alpha} M \\ \bigoplus_ {\lambda \in \mathscr {S}} (1 _ {S _ {\lambda}} \otimes f _ {\lambda}) \Bigg | _ {\downarrow} \qquad \qquad \qquad \qquad \qquad \qquad \Bigg | _ {\varphi (\boldsymbol {f})} \\ \bigoplus_ {\lambda \in \mathscr {S}} (S _ {\lambda} \otimes_ {D _ {\lambda}} V _ {\lambda} ^ {\prime}) \xrightarrow {\alpha^ {\prime}} M ^ {\prime}. \end{array}
\]

这样定义的映射 \(\varphi : \prod_{\lambda \in \mathcal{S}} \mathrm{Hom}_{\mathrm{D}_{\lambda}}(\mathrm{V}_{\lambda}, \mathrm{V}_{\lambda}') \to \mathrm{Hom}_{\mathrm{A}}(\mathrm{M}, \mathrm{M}')\) 是群同构。当我们有 \(\mathrm{M} = \mathrm{M}'\)、\(\mathrm{V}_{\lambda} = \mathrm{V}_{\lambda}'\)（对每个 \(\lambda \in \mathcal{S}\)）、且 \(\alpha = \alpha'\) 时，映射 \(\varphi\) 是从 \(\prod_{\lambda \in \mathcal{S}} \mathrm{End}_{\mathrm{D}_{\lambda}}(\mathrm{V}_{\lambda})\) 到 \(\mathrm{End}_{\mathrm{A}}(\mathrm{M})\) 的环同构。

鉴于命题 8, a) 给出的 M 与 M' 的同型分量的描述，这由 VIII, 第 64 页的定理 3 及 VIII, 第 66 页的命题 5, b) 得出。

推论. —— 设 M 是一个半单 A-模，M' 是一个 A-模。映射 \(u \mapsto (\operatorname{Hom}(1_{\mathrm{S}_{\lambda}}, u))_{\lambda \in \mathcal{S}}\) 从 \(\operatorname{Hom}_{\mathrm{A}}(\mathrm{M}, \mathrm{M}')\) 到

\[
\prod_ {\lambda \in \mathcal {S}} \mathrm{Hom} _ {\mathrm{D} _ {\lambda}} (\mathrm{Hom} _ {\mathrm{A}} (\mathrm{S} _ {\lambda}, \mathrm{M}), \mathrm{Hom} _ {\mathrm{A}} (\mathrm{S} _ {\lambda}, \mathrm{M} ^ {\prime}))
\]

是群同构。当 \(\mathrm{M}'\) 等于 \(\mathrm{M}\) 时，它是从环 \(\operatorname{End}_{\mathrm{A}}(\mathrm{M})\) 到环 \(\prod_{\lambda \in \mathcal{S}} \operatorname{End}_{\mathrm{D}_{\lambda}}(\operatorname{Hom}_{\mathrm{A}}(\mathrm{S}_{\lambda}, \mathrm{M}))\) 的同构。

这是把命题 9 应用于 M 的典范描述与 M' 的底座的典范描述的一个重新表述。

